\documentclass[10pt, a4paper]{article}

% --- Paquetes Esenciales y Configuración de Página ---
\usepackage[utf8]{inputenc}
\usepackage[spanish, es-nodecimaldot]{babel}
\usepackage{amsmath, amssymb, amsfonts}
\usepackage{geometry}
\usepackage{hyperref}
\usepackage{booktabs}
\usepackage{tabularx}
\usepackage{fancyhdr}
\usepackage{microtype}
\usepackage{etoolbox}

% Márgenes estándar
\geometry{margin=2.0cm}

% Configuración de hipervínculos
\hypersetup{
    colorlinks=true,
    linkcolor=black,
    citecolor=black,
    urlcolor=black
}


% Encabezados y pie de página limpios y académicos
\pagestyle{fancy}
\fancyhf{}
\renewcommand{\headrulewidth}{0.4pt}
\fancyhead[L]{\small\itshape Universidad Externado de Colombia}
\fancyhead[R]{\small\bfseries Métodos Numéricos --- Proyecto Integrador (Fase 1)}
\fancyfoot[C]{\thepage}

\begin{document}

% ==============================================================================
% PÁGINA 1: PORTADA
% ==============================================================================
\begin{titlepage}
    \centering
    \vspace*{2.0cm}
    
    {\LARGE\bfseries Proyecto Integrador: Primera Entrega}\\[0.5cm]
    {\large\itshape Derivative-free optimal multiple-root finding family of iterative methods}\\[0.6cm]
    {\Large\textbf{Métodos Numéricos}}\\[2.5cm]
    
    \begin{minipage}{0.85\textwidth}
        \centering
        \textbf{\large Integrantes:}\\[0.3cm]
        Ailyn Sofía Gómez Rodríguez\\[0.15cm]
        Julián Camilo Jiménez\\[0.15cm]
        Julián Steven Duarte\\[1.4cm]
        
        \textbf{\large Docente:}\\[0.3cm]
        Edwin Fernando Muñoz Chipatecua
    \end{minipage}
    
    \vfill
    
    {\normalsize Pregrado en Ciencia de Datos\\
    Departamento de Matemáticas\\
    Universidad Externado de Colombia\\
    Bogotá D.C., Colombia --- Septiembre de 2026}
    \vspace*{0.5cm}
\end{titlepage}

\newpage
\pagenumbering{arabic}
\setcounter{page}{2}

% ==============================================================================
% 1. COMPRENSIÓN DEL ARTÍCULO CIENTÍFICO
% ==============================================================================
\section{Comprensión del Artículo Científico}

\subsection{El problema fundamental: raíces múltiples y degradación de métodos clásicos}
En computación científica y Ciencia de Datos, al resolver ecuaciones no lineales $g(x) = 0$ surgen frecuentemente raíces con \textbf{multiplicidad} $m \ge 2$, donde la función y sus primeras $m-1$ derivadas se anulan simultáneamente. Geométricamente, la gráfica se aplana en un entorno amplio de la raíz, lo que degrada la convergencia cuadrática del método clásico de Newton-Raphson a puramente lineal con factor asintótico $(m-1)/m$. El método de Newton modificado de Schröder restaura el orden cuadrático, pero acarrea la exigencia ineludible de evaluar la derivada analítica $g'(x_k)$ en cada paso. En Ciencia de Datos, donde las funciones provienen de modelos de caja negra o simuladores complejos, la evaluación de derivadas analíticas es inviable, y reemplazarla por diferencias finitas elementales induce errores de truncamiento y cancelaciones numéricas severas \cite{Turner2018}.

\subsection{La propuesta del artículo: esquema óptimo de orden cuatro sin derivadas}
Para superar este dilema, Jerezano, Chicharro y Garrido-Saez (2026) \cite{Jerezano2026} proponen una familia de métodos iterativos en dos pasos que prescinde por completo de derivadas analíticas y alcanza \textbf{orden de convergencia cuatro ($p=4$)}. 

El esquema introduce un desplazamiento cuadrático adaptativo $w_k = x_k + \beta [g(x_k)]^2$ ($\beta \in \mathbb{C} \setminus \{0\}$), aproximando la derivada mediante la diferencia dividida secante $g[w_k, x_k] = \frac{g(w_k)-g(x_k)}{w_k-x_k}$. Con dicha pendiente se genera un predictor intermedio $y_k$, sobre el cual se evalúa la variable de control adimensional $t_k = (g(y_k)/g(x_k))^{1/m}$. Mediante el diseño de una función de peso $H(t)$ que satisface $H(0)=1$, $H'(0)=1$, $H''(0)=4$ y $|H'''(0)| < \infty$, el paso corrector cancela los errores de órdenes inferiores y eleva la convergencia a orden 4.

\textbf{Optimalidad de Kung-Traub:}
La conjetura formulada por Kung y Traub (1974) \cite{KungTraub1974} establece que un método iterativo sin memoria que utiliza $d$ evaluaciones funcionales por iteración puede alcanzar como máximo un orden $p_{\text{máx}} = 2^{d-1}$. El esquema de Jerezano et al. evalúa tres funciones por iteración ($g(x_k), g(w_k), g(y_k)$) y ninguna derivada ($d=3$). Su cota máxima teórica es $2^{3-1} = 4$. Al alcanzar exactamente dicho orden, la familia se clasifica formalmente como \textbf{óptima}.

\subsection{Aritmética de punto flotante y absorción numérica}
Un aspecto numérico crucial radica en la física de punto flotante subyacente a la perturbación secante. En precisión doble estándar IEEE-754 (\texttt{float64}, $\varepsilon_M \approx 2.22 \times 10^{-16}$), la distancia absoluta entre $w_k$ y $x_k$ satisface:
\begin{equation}
    |w_k - x_k| = |\beta| [g(x_k)]^2 \sim \mathcal{O}(e_k^{2m}). \label{eq:absorcion}
\end{equation}
La absorción numérica ocurre cuando la perturbación cae por debajo de la resolución de la mantisa, es decir, cuando $\beta [g(x_k)]^2 < \frac{1}{2} \varepsilon_M |x_k|$. Para la función $g_1(x)$ con $m=2$, $\beta=0.01$ y $|x_k|\approx 1.75$, el estancamiento sobreviene cuando:
\begin{equation}
    |g_1(x_k)| < \sqrt{\frac{\varepsilon_M |x_k|}{2\beta}} \sim 1.4 \times 10^{-7}.
\end{equation}
Dado que cerca de la raíz $|g_1(x)|\approx0.03\,e^2$, el umbral corresponde a $e\approx2.15\times10^{-3}$. Ningún iterado dentro de esa zona puede usarse como punto de partida de un nuevo paso ($0/0$); la precisión final alcanzable depende del último salto de orden 4 y se medirá empíricamente en el experimento.

Por esta razón, los autores en \cite{Jerezano2026} recurrieron a aritmética de precisión variable en MATLAB con \textbf{4000 dígitos de mantisa} para alcanzar la cota $|x_{k+1}-x_k| + |g(x_k)| < 10^{-200}$. En nuestro proyecto en Python, emplearemos la biblioteca \texttt{mpmath} \cite{mpmath2024} con alta precisión para replicar estos resultados y documentar el estancamiento en \texttt{float64}.

\subsection{Resultados numéricos reportados por los autores}
En las Tablas 1, 2 y 3 de \cite{Jerezano2026}, los autores compararon las familias propuestas $M_1$ (acelerador polinomial) y $M_2$ (acelerador racional) frente a métodos de la literatura: Sharma y Sharma (SS) \cite{Sharma2010}, Behl, Cordero, Motsa y Torregrosa (BH) \cite{Behl2016}, Kumar et al. (KK) \cite{Kumar2020b} y Behl et al. (BB) \cite{Behl2021a}:
\begin{enumerate}
    \item \textbf{Orden ACOC:} En todos los casos convergentes se confirmó empíricamente un ACOC de $4.0000$ en 5 a 8 iteraciones.
    \item \textbf{Fallas en la literatura:} El método BB no convergió en ninguno de los tres problemas evaluados (\textit{nc}). Asimismo, los métodos SS y BH no lograron converger en el problema $g_3(x)$.
    \item \textbf{Comparación entre $M_2$ y KK:} El método KK converge en todos los problemas analizados. Para $g_1(x)$ con $x_0 = 1.9$, KK converge en 6 iteraciones mientras que $M_2$ requiere 7 iteraciones. Aunque el texto de \cite{Jerezano2026} afirma que $M_2$ es el más rápido en $g_1(x)$, sus datos en la Tabla 1 muestran que KK requirió una iteración menos. El residuo menor de $M_2$ ($1.77 \times 10^{-1796}$ frente a $1.85 \times 10^{-1661}$ de KK) probablemente se debe a la iteración adicional efectuada y no a una velocidad intrínseca superior. En $g_3(x)$ con $x_0 = 3.4$, $M_1$ y $M_2$ alcanzan la solución en 6 iteraciones frente a 7 de KK.
\end{enumerate}

\subsection{Supuestos y delimitación de limitaciones}
Diferenciamos las limitaciones explícitas del artículo de nuestras inferencias analíticas:

\textbf{Limitaciones declaradas por los autores en el artículo \cite{Jerezano2026}:}
\begin{itemize}
    \item \textbf{Dominio del análisis de estabilidad:} El estudio de dinámica compleja se restringe al polinomio cuadrático $p(z) = (z-1)^2$. Se demostró que $M_2$ carece por completo de puntos fijos extraños para dicho polinomio, mientras que $M_1$ posee puntos fijos extraños pero todos son repulsores (lo que previene converger a raíces falsas). La extensión multidimensional se declaró como trabajo futuro.
    \item \textbf{Comportamiento respecto al parámetro $\beta$:} Al crecer $|\beta|$, las cuencas de atracción de la raíz sufren un encogimiento continuo (Figuras 3 a 9 del artículo). Por ello se adopta $\beta = 1/100$ en las pruebas numéricas.
\end{itemize}

\textbf{Limitaciones e inferencias identificadas por el equipo de trabajo:}
\begin{itemize}
    \item \textbf{Multiplicidad errónea $m' \neq m$:} Introducir una multiplicidad estimada $m'$ distinta de la real provoca que la convergencia se degrade a lineal con un factor que depende de $H$ y de $m'/m$, que se medirá empíricamente en la experimentación.
    \item \textbf{Hipótesis sobre los fallos de SS y BH en $g_3$:} El artículo únicamente registra \textit{nc} para SS y BH en $g_3(x)$. Como hipótesis, los fallos podrían deberse a la sensibilidad a la semilla inicial o a la captura por raíces vecinas (las raíces en $4$ y $5$ están cerca de $x_0 = 3.4$); ambas opciones se contrastarán empíricamente.
    \item \textbf{Rama principal en potencias fraccionarias:} La variable $t_k = (g(y_k)/g(x_k))^{1/m}$ puede arrojar cocientes negativos o complejos. Se adopta la rama principal del logaritmo ($\operatorname{Log}(z) = \ln|z| + i \operatorname{Arg}(z)$ con $-\pi < \operatorname{Arg}(z) \le \pi$) para asegurar una definición determinista de $t_k$.
\end{itemize}

% ==============================================================================
% 2. FORMULACIÓN DEL PROBLEMA COMPUTACIONAL
% ==============================================================================
\section{Formulación del Problema Computacional}

\subsection{Contexto en Ciencia de Datos y objetivo del proyecto}
En Ciencia de Datos, la calibración de modelos y la optimización de pérdidas complejas requieren localizar puntos estacionarios donde el gradiente se anula ($\nabla L(\theta) = 0$). Cuando las funciones de pérdida son evaluadas mediante simulaciones de caja negra y los mínimos presentan curvaturas casi nulas (análogos a raíces múltiples), los solvers estándar se ralentizan o fallan. El objetivo del proyecto consiste en implementar en Python, de forma estructurada y reproducible en el repositorio oficial (\url{https://github.com/ailynsof/Proyecto-M-todos-Numericos}), las familias óptimas de orden 4 de Jerezano et al. (2026), verificando su convergencia y comparándolas con los esquemas del curso.

\subsection{Formulación matemática del algoritmo predictor-corrector}
El esquema iterativo en dos pasos se define mediante:
\begin{align}
    w_k &= x_k + \beta [g(x_k)]^2, \quad \beta \in \mathbb{C} \setminus \{0\}, \label{eq:wk} \\[1pt]
    g[w_k, x_k] &= \frac{g(w_k) - g(x_k)}{w_k - x_k}, \label{eq:dd} \\[1pt]
    y_k &= x_k - m \frac{g(x_k)}{g[w_k, x_k]}, \label{eq:yk} \\[1pt]
    t_k &= \exp\left( \frac{1}{m} \operatorname{Log}\left(\frac{g(y_k)}{g(x_k)}\right) \right), \label{eq:tk} \\[1pt]
    x_{k+1} &= x_k - m \, H(t_k) \frac{g(x_k)}{g[w_k, x_k]}. \label{eq:xk1}
\end{align}
Se implementarán dos familias representativas:
\begin{enumerate}
    \item \textbf{Método $M_1$ (Acelerador Polinomial):} $H_1(t) = 1 + t + 2t^2$.
    \item \textbf{Método $M_2$ (Acelerador Racional):} $H_2(t) = \frac{1 - t}{1 - 2t}$.
\end{enumerate}



% ==============================================================================
% 3. OBTENCIÓN DE LA INFORMACIÓN Y PROBLEMAS DE PRUEBA
% ==============================================================================
\section{Obtención de la Información y Problemas de Prueba}

\subsection{Naturaleza determinista de la información requerida}
En estricto cumplimiento de la guía, se establece que \textbf{no se utilizarán series con ruido estocástico}. El proyecto utilizará \textbf{funciones matemáticas analíticas deterministas} de prueba (\textit{benchmark}) procedentes de problemas de termodinámica, física térmica y álgebra lineal documentados en \cite{Jerezano2026}. El conocimiento del valor analítico exacto de la raíz $\alpha$ permite evaluar con certeza el error absoluto $|x_k - \alpha|$, el residuo $|g(x_k)|$ y la tasa empírica ACOC sin perturbaciones ajenas al algoritmo.

\subsection{Descripción de las funciones benchmark seleccionadas}
Se implementarán los tres problemas de prueba propuestos por los autores:
\begin{enumerate}
    \item \textbf{Problema 1 (Ecuación de Van der Waals en Termodinámica):}
    \begin{equation}
        g_1(x) = x^3 - 5.22 x^2 + 9.0825 x - 5.2675 = (x - 1.75)^2 (x - 1.72). \label{eq:g1}
    \end{equation}
    Modela el comportamiento de gases reales. Presenta una raíz doble ($m=2$) en $\alpha = 1.75$ y una raíz simple vecina en $1.72$. La cercanía extrema ($\Delta = 0.03$) desafía la estabilidad frente a la atracción hacia la raíz simple.
    \item \textbf{Problema 2 (Ley de Radiación Térmica de Planck):}
    \begin{equation}
        g_2(x) = \left( e^{-x} - 1 + \frac{x}{5} \right)^3. \label{eq:g2}
    \end{equation}
    Determina la longitud de onda de máxima densidad espectral en la radiación de cuerpo negro. Posee una raíz triple ($m=3$) en $\alpha \approx 4.9651142317$. Evalúa la robustez ante funciones trascendentes no polinomiales.
    \item \textbf{Problema 3 (Polinomio Característico en Álgebra Lineal):}
    \begin{equation}
        g_3(x) = x^9 - 29x^8 + 349x^7 - 2261x^6 + 8455x^5 - 17663x^4 + 15927x^3 + 6993x^2 - 24732x + 12960. \label{eq:g3}
    \end{equation}
    Ecuación característica de una matriz cuadrada $9 \times 9$, con raíz de multiplicidad $m=4$ en $\alpha = 3$. Se evaluará en su forma polinomial expandida \eqref{eq:g3} tal como en el artículo base para estudiar los efectos de acumulación de redondeo numérico frente a su representación factorizada $(x-3)^4 (x-8)(x-5)(x-4)(x-1)(x+1)$.
\end{enumerate}

El Cuadro~\ref{tab:problemas} resume las características y semillas iniciales de las funciones de prueba. Nótese que para $g_3(x)$ la semilla es $x_0 = 3.4$ en estricta coincidencia con la Tabla 3 de \cite{Jerezano2026}:

\begin{table}[htbp]
    \centering
    \small
    \caption{Características y parámetros de las funciones de prueba del proyecto.}
    \label{tab:problemas}
    \begin{tabularx}{\textwidth}{p{4.2cm} c c c c X}
        \toprule
        \textbf{Función de prueba $g(x)$} & \textbf{Contexto} & \textbf{Raíz $\alpha$} & \textbf{Mult. ($m$)} & \textbf{Semillas ($x_0$)} & \textbf{Desafío numérico principal} \\
        \midrule
        $g_1(x) = (x-1.75)^2(x-1.72)$ & Termodinámica & $1.75$ & 2 & $1.9, \, 3.0$ & Proximidad crítica a raíz simple ($\Delta=0.03$). \\
        $g_2(x) = (e^{-x} - 1 + x/5)^3$ & Física Térmica & $\approx 4.96511$ & 3 & $5.3, \, 10.0$ & Presencia de términos exponenciales trascendentes. \\
        $g_3(x) = \text{Expandido \eqref{eq:g3}}$ & Álgebra Lineal & $3.0$ & 4 & $2.8, \, 3.4$ & Multiplicidad 4 y coeficientes de gran magnitud. \\
        \bottomrule
    \end{tabularx}
\end{table}

% ==============================================================================
% 4. PLAN DE IMPLEMENTACIÓN, CRITERIOS DE COMPARACIÓN Y AVANCES
% ==============================================================================
\section{Plan de Implementación, Criterios y Avances}

\subsection{Métodos numéricos a programar en Python}
En el módulo \texttt{src/} del repositorio del proyecto se desarrollarán los siguientes métodos:
\begin{enumerate}
    \item \textbf{Familias del artículo base:} Métodos $M_1$ y $M_2$ con $\beta = 0.01$. Se programará un barrido complementario con $\beta \in \{0.001, 0.01, 1.0\}$ para evaluar empíricamente el radio de convergencia (aclarando que el análisis formal de dinámica compleja en el plano de Julia queda fuera del alcance de la asignatura).
    \item \textbf{Métodos de contraste del curso y de la literatura:}
    \begin{itemize}
        \item \textit{Newton modificado de Schröder ($p=2$):} Método clásico con derivada analítica $g'(x_k)$.
        \item \textit{Steffensen modificado ($p=2$):} Método adaptativo sin derivadas basado en $w_k = x_k + \beta g(x_k)$ \cite{Traub1964}.
        \item \textit{Kumar et al. (KK, $p=4$) \cite{Kumar2020b}:} El competidor óptimo libre de derivadas de orden 4 más cercano reportado en la literatura, definido formalmente como:
        \begin{equation}
            y_k = x_k - m \frac{g(x_k)}{g[w_k, x_k]}, \qquad x_{k+1} = y_k - \frac{(m+2)t_k}{1-2t_k} \frac{g(x_k)}{g[w_k, x_k] + 2g[w_k, y_k]}, \label{eq:kk}
        \end{equation}
        donde $w_k = x_k + \beta g(x_k)$ y $t_k = (g(y_k)/g(x_k))^{1/m}$.
    \end{itemize}
\end{enumerate}

\subsection{Criterios técnicos de evaluación computacional}
Para evaluar y contrastar cuantitativamente el desempeño de los métodos, se establecen los siguientes criterios:
\begin{enumerate}
    \item \textbf{Criterio de Parada y Tolerancia Numérica:} Se adoptará la cota $|x_{k+1} - x_k| + |g(x_k)| < \text{TOL}$ con un tope de iteraciones $k_{\text{máx}} = 100$.
    \item \textbf{Comportamiento en Flotante Standard (\texttt{float64}):} En \texttt{float64} el método se estanca cuando $w_k = x_k$, es decir, cuando $|g(x_k)| \lesssim \sqrt{\varepsilon_M |x_k| / (2\beta)} \approx 1.39 \times 10^{-7}$ para $g_1$. Se implementará una guarda condicional (\texttt{if w == x: break}) y se reportará el nivel de estancamiento observado como resultado empírico del experimento.
    \item \textbf{Aritmética de Alta Precisión con \texttt{mpmath}:} Para reproducir las tolerancias del artículo ($\text{TOL} = 10^{-200}$) y alcanzar residuos extremos sin absorción numérica, se utilizará la biblioteca \texttt{mpmath} fijando la mantisa en 4000 dígitos decimales \cite{mpmath2024}.
    \item \textbf{Orden Computacional Empírico (ACOC):} Calcularemos el Orden de Convergencia Computacional Aproximado formulado por Cordero y Torregrosa \cite{Cordero2007}:
    \begin{equation}
        \text{ACOC} \approx \frac{\ln |(x_{k+1} - x_k) / (x_k - x_{k-1})|}{\ln |(x_k - x_{k-1}) / (x_{k-1} - x_{k-2})|}. \label{eq:acoc}
    \end{equation}
    Un valor cercano a $4.0000$ confirmará en la práctica la tasa de convergencia cuártica.
    \item \textbf{Sensibilidad a la Semilla Inicial ($x_0$) y Estabilidad:} Se evaluará el radio de atracción probando semillas cercanas y lejanas (ej. $x_0 = 1.9$ y $x_0 = 3.0$ en $g_1$), verificando la presencia de bifurcaciones, oscilaciones o captura por raíces vecinas.
    \item \textbf{Eficiencia Teórica y Tiempo de CPU:} Se contrastará el índice de Kung-Traub ($EI = p^{1/d}$), donde las familias propuestas ($EI = 4^{1/3} \approx 1.587$) superan a Newton modificado ($EI = 2^{1/2} \approx 1.414$), y se medirá el tiempo de procesador en milisegundos promediando 50 ejecuciones independientes.
\end{enumerate}

\subsection{Avance de implementación computacional preliminar en Python}
Como evidencia de avance correspondiente a esta primera entrega, se implementaron las familias $M_1$ y $M_2$ en el script ejecutable \texttt{src/avance\_fase1.py} del repositorio. La ejecución sobre la función $g_1(x)$ con semilla $x_0 = 1.9$ y 4000 dígitos de precisión en \texttt{mpmath} produjo los resultados sintetizados en el Cuadro~\ref{tab:avance}:

\begin{table}[htbp]
    \centering
    \small
    \caption{Resultados preliminares de ejecución en Python (\texttt{src/avance\_fase1.py}) para $g_1(x)$ con $x_0 = 1.9$.}
    \label{tab:avance}
    \begin{tabularx}{\textwidth}{c c c c c X}
        \toprule
        \textbf{Método} & \textbf{Iteraciones ($k$)} & \textbf{$|x_{k+1}-x_k|$} & \textbf{Residuo $|g(x_{k+1})|$} & \textbf{ACOC} & \textbf{Observación de Convergencia} \\
        \midrule
        $M_1$ (Polinomial) & 7 & $9.5190 \times 10^{-537}$ & $0.0$ & $4.0000$ & Cumple tolerancia en $k=7$; coincide con \cite{Jerezano2026}. \\
        $M_2$ (Racional) & 7 & $1.3790 \times 10^{-776}$ & $5.5295 \times 10^{-2012}$ & $4.0000$ & Cumple tolerancia en $k=7$; paso y ACOC concuerdan con \cite{Jerezano2026}. \\
        \bottomrule
    \end{tabularx}
\end{table}

En el Cuadro~\ref{tab:avance}, el paso se define como $|x_{k+1}-x_k|$ y el residuo como $|g(x_{k+1})|$, evaluado tras completar el paso corrector. El número de iteraciones ($k=7$), la magnitud del paso y el ACOC ($4.0000$) concuerdan con los reportados en \cite{Jerezano2026}; el residuo de $M_2$ difiere ($5.53 \times 10^{-2012}$ frente a $1.77 \times 10^{-1796}$ del artículo), discrepancia posiblemente atribuible a particularidades en la definición del punto de evaluación del residuo en el entorno original de MATLAB frente a \texttt{mpmath}, comprometiéndonos a verificarlo en la Fase 2 imprimiendo el error absoluto $e_k = |x_k - \alpha|$ en cada iteración.

% ==============================================================================
% DECLARACIÓN DE USO DE IA
% ==============================================================================
\section*{Declaración sobre el uso de herramientas de Inteligencia Artificial}
\small
En cumplimiento de los lineamientos de integridad académica de la Universidad Externado de Colombia y de la asignatura Métodos Numéricos, los autores declaramos que hemos utilizado herramientas de Inteligencia Artificial generativa como apoyo técnico en la diagramación del documento en \LaTeX{}, la revisión gramatical de estilo, la estructuración preliminar del script exploratorio \texttt{src/avance\_fase1.py} y la discusión analítica de conceptos teóricos y numéricos. Los integrantes del equipo comprendemos integralmente, sustentamos y asumimos la total autoría y responsabilidad del contenido matemático, la formulación computacional y los resultados de este documento.

% ==============================================================================
% REFERENCIAS BIBLIOGRÁFICAS
% ==============================================================================
\begin{thebibliography}{99}
\footnotesize
\setlength{\itemsep}{-0.5pt}
\bibitem{Jerezano2026} Jerezano, J. H., Chicharro, F. I., \& Garrido-Saez, N. (2026). Derivative-free optimal multiple-root finding family of iterative methods. \textit{Computational and Applied Mathematics}, 45:304.
\bibitem{KungTraub1974} Kung, H. T., \& Traub, J. F. (1974). Optimal order of one-point and multipoint iteration. \textit{Journal of the ACM}, 21(4), 643--651.
\bibitem{Traub1964} Traub, J. F. (1964). \textit{Iterative Methods for the Solution of Equations}. Prentice-Hall.
\bibitem{Sharma2010} Sharma, J. R., \& Sharma, R. (2010). Modified Jarratt method for computing multiple roots. \textit{Applied Mathematics and Computation}, 217, 878--881.
\bibitem{Behl2016} Behl, R., Cordero, A., Motsa, S. S., \& Torregrosa, J. R. (2016). An optimal fourth-order family of methods for multiple roots and its dynamics. \textit{Numerical Algorithms}, 71, 775--796.
\bibitem{Kumar2020b} Kumar, S., Kumar, D., Sharma, J. R., Cesarano, C., Agarwal, P., \& Chu, Y. M. (2020). An optimal fourth order derivative-free numerical algorithm for multiple roots. \textit{Symmetry}, 12, 1038.
\bibitem{Behl2021a} Behl, R., Bhalla, S., Magreñán, Á. A., \& Moysi, A. (2021). An optimal derivative free family of Chebyshev--Halley's method for multiple zeros. \textit{Mathematics}, 9, 546.
\bibitem{Cordero2007} Cordero, A., \& Torregrosa, J. R. (2007). Variants of Newton’s method using fifth-order quadrature formulas. \textit{Applied Mathematics and Computation}, 190, 686--698.
\bibitem{Turner2018} Turner, P. R., Arildsen, T., \& Kavanagh, K. (2018). \textit{Applied Scientific Computing with Python}. Springer-Verlag.
\bibitem{mpmath2024} Johansson, F. et al. (2024). \textit{mpmath: Python library for arbitrary-precision floating-point arithmetic} (Version 1.4.1). \url{https://mpmath.org/}.
\end{thebibliography}

\end{document}
